\subsection{An elementary proof of the arithmetic centralizer assertion}

\paragraph{Statement B1 (Proved: R3).}\label{res:B1}
Fix an algebraic closure $\overline{\mathbf Q}$, with its profinite absolute
Galois group $G_{\mathbf Q}$. For every number field
$k\subset\overline{\mathbf Q}$, put $F=k(\mu_\infty)$. Then
\[
 C_{G_{\mathbf Q}}(G_F)=1.
\]
Here $G_F=\operatorname{Gal}(\overline{\mathbf Q}/F)$ is a closed subgroup,
and the centralizer is taken in the full profinite group, without passing to
an abelian or pro-$\ell$ quotient. This proves the starred assertion in
cote 149, \S2, manuscript p.~22, scan 45.

\begin{proof}
Replacing $k$ by its finite Galois closure $k_0/\mathbf Q$ gives the inclusion
$G_{k_0(\mu_\infty)}\subseteq G_F$, and hence the inclusion of centralizers
in the opposite direction. It therefore suffices to consider $k/\mathbf Q$
Galois. Fix $\sigma\in C_{G_{\mathbf Q}}(G_F)$.

For $a\in F$ and $h\in G_F$ we have
$h(\sigma a)=\sigma(ha)=\sigma a$. Applying the same argument to
$\sigma^{-1}$ proves $\sigma(F)=F$. For arbitrary
$x\in\overline{\mathbf Q}$ the identities
\[
 h(x)=x\quad\Longleftrightarrow\quad h(\sigma x)=\sigma x
 \qquad(h\in G_F)
\]
identify the two closed stabilizers in $G_F$. By the Galois correspondence,
\begin{equation}\label{eq:B-stabilizer}
 F(x)=F(\sigma x).
\end{equation}

We record explicitly the choice of split prime needed below. A nonconstant
polynomial $f\in\mathbf Z[T]$ has values divisible by infinitely many
rational primes. Indeed, after translating $T$ we have $f(0)=c\ne0$.
For any finite set $P$ of primes, enlarge $P$ to contain the prime divisors
of $c$ and put $b=\prod_{p\in P}p$. Then
$f(cbt)/c$ is an integer congruent to $1$ modulo every prime in $P$;
for some integer $t$ its absolute value exceeds $1$. It has a prime divisor
outside $P$. Apply this to the minimal polynomial of an integral primitive
element of a finite Galois number field $k'/\mathbf Q$. Outside the finitely
many primes dividing its discriminant or the index of its generated order,
a root modulo $p$ gives a residue-degree-one prime of $k'$. All residue
degrees are equal in an unramified Galois extension; thus $p$ splits
completely. This proves the existence of completely split primes outside
any prescribed finite set, without an appeal to a density theorem.

Suppose that $\sigma$ acts nontrivially on a finite Galois subfield
$k'\subset F$. Choose a rational prime $p$ which splits completely in $k'$
and is unramified in $k$. The action of
$\operatorname{Gal}(k'/\mathbf Q)$ on the primes above $p$ is simply
transitive, so $\sigma$ moves each such prime. Choose one, $\mathfrak p$.
The Chinese remainder theorem, applied modulo $\mathfrak p^2$ and the
other primes above $p$, gives $a\in k'^*$ such that
\[
 v_{\mathfrak p}(a)=1,\qquad
 v_{\mathfrak q}(a)=0\quad(\mathfrak q\mid p,\ \mathfrak q\ne\mathfrak p).
\]
The normalization here is $v(p)=1$. Extend $v_{\mathfrak p}$ to $F$.
Its value group is exactly
\begin{equation}\label{eq:B-value-group}
 v(F^*)=\frac1{p-1}\mathbf Z[1/p].
\end{equation}
To check the equality, roots of unity of order prime to $p$ belong to the
maximal unramified extension of $\mathbf Q_p$: they lift the roots in finite
residue fields by the simple-root form of Hensel's lemma. The extensions
$\mathbf Q_p(\mu_{p^r})/\mathbf Q_p$ are totally ramified of degree
$(p-1)p^{r-1}$; this degree and total ramification follow from the
Eisenstein polynomial $\Phi_{p^r}(1+T)$.
Since $k/\mathbf Q$ is unramified at $p$, its completions do not add a
ramification factor. The value groups of the finite composita are therefore
contained in, and the $p$-power cyclotomic fields already realize, the union
in \eqref{eq:B-value-group}. In particular this argument requires
unramifiedness in $k$, not merely complete splitting in $k'$.

Choose a prime $\ell\nmid p(p-1)$. Equation \eqref{eq:B-value-group} gives
$1\notin\ell v(F^*)$, so $a\notin(F^*)^\ell$. Choose
$\alpha^\ell=a$. The extension $F(\alpha)/F$ is cyclic of degree $\ell$,
because $F$ contains $\mu_\ell$. Equation \eqref{eq:B-stabilizer} gives
$F(\alpha)=F(\sigma\alpha)$. If a generator $\tau$ sends
$\alpha$ to $\zeta_\ell\alpha$, then
$\tau(\sigma\alpha)/(\sigma\alpha)\in\mu_\ell$.
The $\zeta_\ell^j$-eigenspace for $\tau$ in
$F(\alpha)=\bigoplus_{i=0}^{\ell-1}F\alpha^i$ is $F\alpha^j$.
Consequently
\[
 \sigma\alpha=c\alpha^j,\qquad
 \sigma(a)=c^\ell a^j,
 \qquad c\in F^*,\quad 1\le j\le\ell-1.
\]
The condition $j\ne0\bmod\ell$ follows from $\sigma\alpha\notin F$.
At our chosen valuation, $v(\sigma a)=0$ whereas $v(a)=1$. Thus
$0=\ell v(c)+j$. Since
$\ell v(F^*)\cap\mathbf Z=\ell\mathbf Z$, this contradicts
$1\le j\le\ell-1$.

We have proved that $\sigma$ fixes every finite Galois subfield of $F$;
their union is $F$, since $F/\mathbf Q$ is Galois. Finally, for an arbitrary
finite Galois number field $k''$ containing $k$, the inclusion
$G_{k''(\mu_\infty)}\subset G_F$ shows that $\sigma$ centralizes that
subgroup too. The argument just given, with $k''$ in place of $k$, proves
that $\sigma$ fixes $k''(\mu_\infty)$ pointwise. These fields exhaust
$\overline{\mathbf Q}$, so $\sigma=1$.
\end{proof}



\paragraph{Audit of the supplied R3 argument.}
Every implication in its proposed chain is valid after the preceding two
additions: the prime $p$ is unramified in the original finite coefficient
field, and the chosen Kummer class is nonzero. In the final passage one may
use finite Galois number fields $k''$ containing $k$; it is unnecessary to
assert that an arbitrary finite extension of $F$ is Galois over
$\mathbf Q$. It is true that every finite $F'/F$ has the form
$k_1(\mu_\infty)$ for a number field $k_1$: choose finitely many algebraic
generators of $F'/F$ and put them together with $k$ into $k_1$.

\subsection{A general Hilbertian criterion and the fields in Part B(c)}

\paragraph{Statement B2 (Proved: a stronger arithmetic criterion).}\label{res:B2}
For a Hilbertian field $k$ and a separable closure $k^{\rm sep}$, every
nontrivial closed normal subgroup $H\triangleleft G_k$ satisfies
$C_{G_k}(H)=1$. Consequently, for every algebraic Galois extension
$M/k$ in $k^{\rm sep}$,
\[
 C_{G_k}(G_M)=
 \begin{cases}
 1,&M\ne k^{\rm sep},\\
 G_k,&M=k^{\rm sep}.
 \end{cases}
\]
The comparison is in the category of profinite groups with continuous
homomorphisms.

\begin{proof}
We use two precisely identified inputs from the opened primary texts.
Jarden \cite[Proposition 2.5]{B_Jarden2011} treats a Hilbertian field $k$
and an algebraic Galois extension $N/k$ that is not separably closed. Its
conclusion is that $G_N$ is almost locally free: for each integer $e\ge1$,
almost every $e$-tuple, with respect to Haar measure, generates a free
profinite group of rank $e$ (definition (2) in the introduction).
Haran--Jarden \cite[Corollary 2.5]{B_HaranJarden1991} states that the
separable closure of a Hilbertian field cannot be the compositum of two
proper Galois extensions of the field. Their Theorem 2.4 gives the stronger
input: the compositum of two Galois extensions, neither contained in the
other, is Hilbertian. These statements include infinite Galois extensions.

There is no nontrivial closed abelian normal subgroup $A$ of $G_k$.
Otherwise its fixed field is a proper Galois extension of $k$, and Jarden's
Proposition 2.5 supplies a pair in $A$ generating a free profinite group of
rank two. Such a group has the nonabelian quotient $S_3$, whereas every
subgroup and every quotient of the abelian group $A$ is abelian, a
contradiction.

Now put $C=C_{G_k}(H)$. Centralizers are closed, since each commutation
condition is closed, and $C$ is normal since $H$ is normal. The subgroup
$H\cap C=Z(H)$ is closed, abelian, and normal in $G_k$, so it is trivial.
If $C$ were nontrivial, the fixed fields
$(k^{\rm sep})^H$ and $(k^{\rm sep})^C$ would both be proper Galois
extensions of $k$. The identity
\[
 G_{(k^{\rm sep})^H(k^{\rm sep})^C}=H\cap C=1
\]
would make their compositum $k^{\rm sep}$, contrary to the cited
Corollary 2.5. Hence $C=1$. The second displayed formula follows by taking
$H=G_M$; when $H=1$, every element centralizes it.
\end{proof}

This is a derivation from the two opened theorems, not a claim of priority.
Our search included the phrases ``absolute Galois centralizer normal
subgroup'', ``Hilbertian prosolvable normal subgroup'', and ``centralizer
Galois abelian extension''. We did not find the exact formula of Statement
B1 stated separately in the opened papers. Statement B2 supplies a
literature-based proof of it after taking a Galois closure of $k$; Statement
B1 independently checks the proposed elementary proof.

\paragraph{Statement B3 (Proved: all the named fields).}\label{res:B3}
For each field in the following table the indicated centralizer is trivial.
Here $k$ is any number field; for the first two rows a finite Galois closure
of $k/\mathbf Q$ may be used when $k(\mu_{\ell^\infty})$ or
$k(\mu_\infty)$ itself is not Galois over $\mathbf Q$.
\begin{center}
\begin{tabular}{ll}
\hline
$M$ & Centralizer in $G_{\mathbf Q}$ of $G_M$\\\hline
$k(\mu_{\ell^\infty})$ & $1$\\
$k(\mu_\infty)$ & $1$\\
$\mathbf Q^{\rm ab}$ & $1$\\
the maximal nilpotent extension of $\mathbf Q$ & $1$\\
the maximal solvable extension of $\mathbf Q$ & $1$\\
the field $\mathbf Q^{\rm tr}$ of totally real algebraic numbers & $1$\\\hline
\end{tabular}
\end{center}

\begin{proof}
The first row follows directly from
$G_{k(\mu_\infty)}\subset G_{k(\mu_{\ell^\infty})}$ and Statement B1.
The second is Statement B1. For all remaining rows we use Statement B2
with $k=\mathbf Q$, which is Hilbertian by Hilbert's irreducibility theorem
(opened as \cite[\S4, Theorem 1]{B_Geyer2012}). We must check that each
Galois extension is proper.

The splitting field of $T^3-2$ has group $S_3$: Eisenstein's criterion at
$2$ gives irreducibility, and its discriminant $-108$ is not a square.
It cannot be contained in an abelian extension or in a pro-nilpotent
extension of $\mathbf Q$. This proves properness for $\mathbf Q^{\rm ab}$
and the maximal nilpotent extension, without using the Kronecker--Weber
identification. The splitting field of $T^5-T-1$ has group $S_5$.
Here is a finite certificate with the inference spelled out. Its reduction
modulo $3$ is irreducible, so its rational Galois group is transitive of
degree five. Its squarefree reduction modulo $2$ has irreducible factors
of degrees two and three, so the group has an element of cycle type
$(2)(3)$ and hence a transposition. A transitive group of prime degree is
primitive. The graph whose edges are the conjugates of that transposition
has connected components forming blocks; it is therefore connected.
Transpositions along a connected graph generate $S_5$.
Thus this Galois group is $S_5$, and is not solvable. It cannot be a quotient
of the prosolvable group of the maximal solvable extension. The field
$\mathbf Q^{\rm tr}$ is Galois, because being totally real is invariant
under conjugation, and it omits $i$. It too is proper.

The program \texttt{code/B\_field\_certificates.py} computes the indicated
factorizations, the discriminants, the derived series orders
$120,60,60$ of $S_5$, and the three Sylow-$2$ subgroups of $S_3$.
Its exact output is attached. The factorization-to-cycle inference follows
by reducing the roots at an unramified prime: the residue Frobenius permutes
the roots in cycles whose lengths are the irreducible factor degrees.
\end{proof}

\subsection{Absolute groups of finitely generated fields and inverse limits}

\paragraph{Statement B4 (Proved: the assertion for $E_K$).}\label{res:B4}
Let $K/\mathbf Q$ be a finitely generated field, fix an algebraic closure
$\overline K$, and write $E_K=G_K$. For every open subgroup $E'\subset E_K$
let $E'^\circ$ be the kernel of its action on all roots of unity in
$\overline K$. Then
\[
 C_{E_K}(E'^\circ)=1.
\]
This proves the assertion accompanying the starred clause of scan 45 in
its stated field-theoretic generality.

\begin{proof}
The field $K$ is Hilbertian: the opened source
\cite[\S4, Theorem 2(e)]{B_Geyer2012} states this for every infinite field
finitely generated over its prime field. Choose a finite Galois extension
$K_1/K$ whose absolute Galois group is contained in $E'$ (the normal core
of $E'$ provides it). The extension $K_1(\mu_\infty)/K$ is Galois and
$G_{K_1(\mu_\infty)}\subset E'^\circ$.
It is a proper subfield of $\overline K$. Indeed $K_1$ is Hilbertian and
$K_1(\mu_\infty)/K_1$ is abelian; the opened Kuyk statement in
\cite[\S4, Theorem 2(g)(2)]{B_Geyer2012} says that every algebraic abelian
extension of a Hilbertian field is Hilbertian. A separably closed field is
not Hilbertian: $Y^2-T$ is irreducible in its rational function field but
every nonzero specialization is reducible in characteristic zero.
Statement B2, applied to the nontrivial normal subgroup
$G_{K_1(\mu_\infty)}\triangleleft G_K$, now gives the result.
\end{proof}

\paragraph{Statement B5 (Proved, with the rank-one correction).}\label{res:B5}
The property
\[
 \text{every open subgroup $H\subset G$ has $C_G(H)=1$}
 \tag{B-slim}
\]
is preserved by extensions of profinite groups and by projective limits
whose transition homomorphisms are surjective. A free profinite group of
finite rank at least two has this property. The trivial group has it.
The free profinite group of rank one does not: it is $\widehat{\mathbf Z}$,
and its centralizer of every subgroup is the whole group. Thus the
corollary as worded on scan 15 needs the nonabelian-rank qualification on
every nontrivial free factor.

\begin{proof}
For a free profinite group $F$ of rank $r\ge2$, every finite-index subgroup
$V$ is free profinite of rank $1+[F:V](r-1)$. One can see the comparison
explicitly: the preimage $V_0$ of $V$ in the dense discrete free group
$F_r$ corresponds to a finite covering of its graph with $r$ loops.
Choosing a spanning tree gives a free basis with
$1+[F:V](r-1)$ elements. Its full profinite topology is induced from $F_r$,
because the core in $F_r$ of any finite-index subgroup of $V_0$ still has
finite index. Therefore its profinite completion maps isomorphically onto
$V$, and $V^{\rm ab}\simeq\widehat{\mathbf Z}^{1+[F:V](r-1)}$.

This also proves torsion-freeness without assuming it. A nonidentity
finite-order $z\in F$ has a nonidentity image in a finite quotient $F/N$.
Put $V=N\overline{\langle z\rangle}$. The finite-order image of $z$ in
the torsion-free group $V^{\rm ab}$ is zero, but the map
$V\twoheadrightarrow V/N$, whose target is the nontrivial cyclic group
generated by $zN$, factors through $V^{\rm ab}$. This is a contradiction.

The center of every free profinite group of rank at least two is trivial
by the opened Nakamura \cite[Corollary (1.3.4)]{B_Nakamura1994}; its
hypotheses are the pro-$\mathcal C$ completion of a finite-rank free group
of rank $r\ge2$, with $\mathcal C$ a full class of finite groups. We take
the class of all finite groups. For an open $H\subset F$, choose an open
normal $V\subset H$. Its rank is at least two, so $Z(V)=1$.
The homomorphism $C_F(V)\to F/V$ is injective, because its kernel is
$Z(V)$. Thus $C_F(V)$ is finite; torsion-freeness makes it trivial.
The inclusion $C_F(H)\subset C_F(V)$ proves (B-slim).

For an exact sequence of profinite groups
$1\to A\xrightarrow{i}B\xrightarrow{q}D\to1$, where $A,D$ have
(B-slim), let $H$ be open in $B$ and $z\in C_B(H)$. The quotient map is
open, so $q(H)$ is open in $D$, and $q(z)$ centralizes it. Hence $q(z)=1$.
Then $z\in A$ centralizes $H\cap A$, open in $A$, and so $z=1$.

Finally let $G=\varprojlim_iG_i$, with continuous surjective transition
maps, each $G_i$ satisfying (B-slim). The projections $q_i:G\to G_i$
are surjective. An open $H\subset G$ contains $\ker q_{i_0}$ for some
index $i_0$: finitely many inverse-image neighborhoods contain an open
neighborhood of $1$ in $H$, and directedness combines their indices.
For each $j\ge i_0$ the subgroup $q_j(H)$ is open in $G_j$, since it is
the inverse image of $q_{i_0}(H)$ under $G_j\to G_{i_0}$.
For $z\in C_G(H)$, its coordinate $q_j(z)$ centralizes $q_j(H)$, and is
therefore $1$. These $j$ are cofinal, so $z=1$ in the inverse limit.
\end{proof}

Here is the passage relevant to the manuscript, with its maps specified.
For a regular integral model $Y$ of a function field $L$, and a dense open
$V\subset Y$, restriction of finite \emph{\'etale} covers gives a
continuous surjection $\pi_1(V)\twoheadrightarrow\pi_1(Y)$.
A connected finite \emph{\'etale} cover of the normal integral $Y$ stays
connected on a dense open: its function field is unchanged. This is the
connected-cover criterion for surjectivity of the fundamental-group map.
Every finite separable extension of $L$ is finite \emph{\'etale} on some
dense open of its normalization, after deleting the branch locus. Thus
\[
 G_L\ \xrightarrow{\sim}\ \varprojlim_{V\subset Y}\pi_1(V)
\]
is the continuous isomorphism induced by restriction to the corresponding
maximal extensions unramified on $V$. The same construction after extending
the constant field gives the geometric inverse limit. Whenever the cofinal
model groups are the successive nonabelian free extensions constructed in
\S1 of the manuscript, Statement B5 gives exactly the missing passage
to the limit; it requires surjectivity of the transition maps, and no
injectivity assertion.

There is also a direct verification of the needed property of the
geometric absolute group, independent of that choice of models. Let
$k$ be the algebraic closure of $\mathbf Q$ in $K$, and put
$\pi_K=G_{K\overline k}$. For positive transcendence degree,
$K\overline k$ is a finite extension of a rational function field
$\overline k(t_1,\ldots,t_d)$, so it is Hilbertian by
\cite[\S4, Theorem 2(a),(d)]{B_Geyer2012}. Statement B2 applied to the
normal core of an open subgroup gives (B-slim) for $\pi_K$. In
transcendence degree zero it is trivial. Projection of an element of
$C_{G_K}(E'^\circ)$ to $G_k$ centralizes an open cyclotomic kernel, and
is trivial by B1; the element then lies in $\pi_K$ and centralizes
$E'\cap\pi_K$. This provides the manuscript's original reduction as well.

The three marginal questions on scans 14--15 concern, respectively,
the center of the open arithmetic subgroup, the center of a free
profinite group, and its torsion. In this last proof the first is answered
by B1; the second is answered, with rank at least two, by Nakamura's cited
Corollary (1.3.4); and the third is proved above. All three enter the
manuscript's finite-centralizer argument. The alternative proof B4 uses
the Hilbertian criterion instead and does not require those three
auxiliary steps.

\subsection{Arithmetic and geometric fixed groups}

\paragraph{Statement B6 (Proved: the equivalence used at scan 127).}\label{res:B6}
Let $K/\mathbf Q$ be finitely generated, let
\[
 1\longrightarrow P\longrightarrow E\xrightarrow{p}G_K\longrightarrow1
\]
be an exact sequence of profinite groups, and let
$s:G_{K'}\to E\times_{G_K}G_{K'}$ be a continuous section over a finite
extension $K'/K$. We regard its image as a subgroup of $E$ through the
inclusion of the pullback. For every subgroup $\Omega\subset G_{K'}^\circ$
that is open in $G_{K'}^\circ$,
\[
 C_E(s(\Omega))=C_P(s(\Omega)).
\]
In particular the corresponding two assertions of triviality are
equivalent. This equality applies both to the absolute group of a
function field and to the arithmetic fundamental group of a model.

\begin{proof}
An open subgroup $\Omega$ of $G_{K'}^\circ$ contains
$G_{K''}^\circ$ for a finite extension $K''/K'$.
To construct it, choose an open subgroup $W\subset G_{K'}$ with
$W\cap G_{K'}^\circ\subset\Omega$, and let $K''$ be its fixed field.
Then $G_{K''}^\circ=W\cap G_{K'}^\circ$. For
$z\in C_E(s(\Omega))$, the element $p(z)$ centralizes
$G_{K''}^\circ$ in $G_K$, and is trivial by B4. Thus $z\in P$.
The reverse inclusion of the two centralizers is the inclusion $P\to E$.
Every map used is the displayed continuous quotient or subgroup inclusion;
no comparison of fixed groups under a quotient is inferred in the reverse
direction.
\end{proof}

\paragraph{Statement B12 (Proved: Theorem 1 over cyclotomic fields).}\label{res:B12}
Let $k$ be a number field, $F=k(\mu_\infty)$, and let $M/F$ be a finitely
generated field extension. Every open subgroup $H\subset G_M$ has
$C_{G_M}(H)=1$. Since $\mu_\infty\subset M$, its cyclotomic kernel is
$G_M$ itself. For positive transcendence degree, the geometric absolute
group $G_{M\overline F}$ has the same open-subgroup centralizer property;
in transcendence degree zero it is trivial. This proves Theorem 1 of scan
13, including the geometric assertion, over $F$ and its finitely generated
extensions.

\begin{proof}
The abelian extension $F/k$ is Hilbertian by Kuyk's theorem as opened in
\cite[\S4, Theorem 2(g)(2)]{B_Geyer2012}. For $M/F$ algebraic and finitely
generated, $M/F$ is finite and preserves Hilbertianity by that source's
Theorem 2(d). For positive transcendence degree, a transcendence basis
expresses $M$ as a finite extension of a rational function field over
$F$, which is Hilbertian by Theorem 2(a), followed again by (d).
Choose an open normal subgroup $H_0\subset H$. Its fixed field is finite
over $M$, hence Hilbertian and not separably closed, so $H_0\ne1$.
Statement B2 gives $C_{G_M}(H_0)=1$, and hence $C_{G_M}(H)=1$.
For the geometric group, $M\overline F$ is a finite extension of
$\overline F(t_1,\ldots,t_d)$ when $d>0$, so the identical proof applies.
When $d=0$ it is the algebraically closed field $\overline F$, with trivial
absolute Galois group.
\end{proof}
